% from ega3-1-fr.tex
\oldpage[III]{83}
\begin{env}[1.1.3]
\label{III.1.1.3-fr}
Des complexes précédents, on déduit comme d'ordinaire les \emph{$A$-modules
d'homologie et de cohomologie} (G, I, 2.2)
\begin{align}
\label{III.1.1.3.1-fr}
\tag{1.1.3.1}
H_\bullet(\mathbf f,M)&=H_\bullet(K_\bullet(\mathbf f,M)),\\
\label{III.1.1.3.2-fr}
\tag{1.1.3.2}
H^\bullet(\mathbf f,M)&=H^\bullet(K^\bullet(\mathbf f,M)).
\end{align}
On définit d'ailleurs un $A$-isomorphisme
$K_i(\mathbf f,M)\xrightarrow{\sim}K^{r-i}(\mathbf f,M)$ en faisant
correspondre à toute chaîne
$z=\sum(\mathbf e_{i_1}\wedge\cdots\wedge\mathbf e_{i_k})\otimes
z_{i_1\ldots i_k}$ la cochaîne $g_z$ telle que
$g_z(j_1,\ldots,j_{r-k})=\varepsilon z_{i_1\ldots i_k}$, où
$(j_h)_{1\leq h\leq r-k}$ est la suite strictement croissante complémentaire
de la suite strictement croissante $(i_h)_{1\leq h\leq k}$ dans $[1,r]$ et
$\varepsilon=(-1)^\nu$, $\nu$ étant le nombre d'inversions de la permutation
$i_1,\ldots,i_k,j_1,\ldots,j_{r-k}$ de $[1,r]$. On vérifie que
$g_{dz}=d(g_z)$, ce qui donne un isomorphisme
\begin{equation}
\label{III.1.1.3.3-fr}
\tag{1.1.3.3}
H^i(\mathbf f,M)\xrightarrow{\sim}H_{r-i}(\mathbf f,M)
\qquad\text{pour }0\leq i\leq r.
\end{equation}
Dans ce chapitre, nous aurons surtout à considérer les modules de cohomologie
$H^\bullet(\mathbf f,M)$.

Pour un $\mathbf f$ donné, il est immédiat (G, I, 2.1) que
$M\mapsto H^\bullet(\mathbf f,M)$ est un \emph{foncteur cohomologique}
(T, II, 2.1), de la catégorie des $A$-modules dans celle des $A$-modules
gradués, nul pour les degrés $<0$ et $>r$. En outre, on a
\begin{equation}
\label{III.1.1.3.4-fr}
\tag{1.1.3.4}
H^0(\mathbf f,M)=\operatorname{Hom}_A(A/(\mathbf f),M)
\end{equation}
en désignant par $(\mathbf f)$ l'idéal de $A$ engendré par
$f_1,\ldots,f_r$~; cela résulte aussitôt de
(\hyperref[III.1.1.2.3-fr]{1.1.2.3}), et il est clair que
$H^0(\mathbf f,M)$ s'identifie au sous-module de $M$ \emph{annulé par
$(\mathbf f)$}. De même, on a d'après
(\hyperref[III.1.1.2.3-fr]{1.1.2.3})
\begin{equation}
\label{III.1.1.3.5-fr}
\tag{1.1.3.5}
H^r(\mathbf f,M)=M/\Bigl(\sum_{i=1}^r f_iM\Bigr)
=(A/(\mathbf f))\otimes_A M.
\end{equation}
Nous utiliserons le résultat connu suivant, dont nous rappellerons une
démonstration pour être complet~:
\end{env}

\begin{proposition}[1.1.4]
\label{III.1.1.4-fr}
Soient $A$ un anneau, $\mathbf f=(f_i)_{1\leq i\leq r}$ une famille finie
d'éléments de $A$, $M$ un $A$-module. Si, pour $1\leq i\leq r$, l'homothétie
$z\mapsto f_i z$ dans $M_{i-1}=M/(f_1M+\cdots+f_{i-1}M)$ est injective, on a
$H^i(\mathbf f,M)=0$ pour $i\neq r$.
\end{proposition}

\begin{proof}
Il suffit en effet de prouver que $H_i(\mathbf f,M)=0$ pour tout $i>0$ en
vertu de (\hyperref[III.1.1.3.3-fr]{1.1.3.3}). Raisonnons par récurrence sur
$r$, le cas $r=0$ étant trivial. Posons
$\mathbf f'=(f_i)_{1\leq i\leq r-1}$~; cette famille vérifie les conditions de
l'énoncé, donc, si on pose $L_\bullet=K_\bullet(\mathbf f',M)$, on a
$H_i(L_\bullet)=0$ pour $i>0$ par hypothèse, et
$H_0(L_\bullet)=M_{r-1}$ en vertu de
(\hyperref[III.1.1.3.3-fr]{1.1.3.3}) et
(\hyperref[III.1.1.3.5-fr]{1.1.3.5}). Posons pour abréger
$K_\bullet=K_\bullet(f_r)=K_0\oplus K_1$, avec $K_0=K_1=A$,
$d_1:K_1\to K_0$ étant la multiplication par $f_r$~; on a par définition
(\hyperref[III.1.1.1-fr]{1.1.1})
$K_\bullet(\mathbf f,M)=K_\bullet\otimes_A L_\bullet$. Or, on a le lemme
suivant~:

\begin{lemma}[1.1.4.1]
\label{III.1.1.4.1-fr}
Soit $K_\bullet$ un complexe de chaînes formé de $A$-modules libres, nuls sauf
en dimensions $0$ et $1$. Pour tout complexe de chaînes $L_\bullet$ formé de
$A$-modules, on a une suite exacte
\[
0\longrightarrow H_0(K_\bullet\otimes H_p(L_\bullet))
\longrightarrow H_p(K_\bullet\otimes L_\bullet)
\longrightarrow H_1(K_\bullet\otimes H_{p-1}(L_\bullet))
\longrightarrow0
\]
pour tout indice $p$.
\end{lemma}

